\section{Related Work}
\label{sec:related_work}

Our work operates at the intersection of three active fields of research: model merging, neural network compression via pruning, and test-time adaptation. 

\subsection{Model Merging and Task Arithmetic}
Model merging fusions multiple task-specific neural networks into a single cohesive model without undergoing computationally intensive training or accessing full training datasets. Early approaches like \emph{Model Soups} \cite{wortsman2022model} demonstrated that averaging the weights of multiple models fine-tuned on the same task with different hyperparameters improves accuracy and generalization. This concept was extended to multi-task learning via \emph{Task Arithmetic} \cite{ilharco2023editing}, which defines a task vector by subtracting pre-trained base parameters from fine-tuned task-specific parameters. Adding these task vectors together allows a single base model to perform multiple distinct downstream tasks simultaneously.

However, naive task vector addition often causes destructive interference between tasks due to overlapping and conflicting parameter updates. To resolve this interference, several methods have been introduced: \emph{TIES-Merging} \cite{yadav2023ties} resolves sign conflicts and prunes redundant small-magnitude parameter changes; \emph{Fisher-Merging} \cite{matena2021merging} utilizes the diagonal Fisher Information Matrix to weigh parameter importance. More recently, \emph{AdaMerging} \cite{yang2024adamerging} introduced an adaptive layer-wise coefficient optimization scheme that fine-tunes merging weights at test-time on a tiny, unlabeled calibration dataset. Other works like \emph{RegCalMerge} \cite{daheim2023elastic} introduce calibration and regularizations to avoid transductive overfitting, and \emph{PolyMerge} \cite{sanford2023representing} constrains optimization to smooth low-dimensional polynomial subspaces. However, all these methods assume a fully dense merged model, which represents a major deployment hurdle on resource-constrained devices.

\subsection{Model Compression and Weight Pruning}
Network pruning is a foundational compression paradigm designed to reduce deep neural networks' storage footprint and memory bandwidth requirements for edge deployment \cite{han2015learning}. Pruning techniques can be broadly categorized into unstructured pruning (which zeroes out individual unimportant weights) and structured pruning (which removes entire channels, heads, or layers) \cite{li2016pruning,he2017channel}. Unstructured magnitude pruning \cite{han2015learning} remains a highly popular, robust baseline that sorts weights based on their absolute values and removes those below a target percentile threshold. 

The \emph{Lottery Ticket Hypothesis} \cite{frankle2018lottery,chen2020lottery} demonstrated that dense networks contain sparse, trainable subnetworks capable of achieving matching or superior performance when trained from scratch. However, finding these sparse sub-networks is extremely compute-intensive. In the context of model merging, while methods like TIES-Merging \cite{yadav2023ties} prune individual task vectors post-hoc to resolve sign and parameter conflict, they do not optimize the overall merged network under physical storage constraints. Recent work has explored post-training quantization for merged models \cite{detettori2023pruning}, but joint weight-pruning and model merging remains an open, vital area of exploration for physical edge deployment.

Furthermore, recent advances in neural network compression, especially for large language models, have proposed alternative data-free or calibration-dependent pruning metrics. For example, \emph{Wanda} \cite{sun2023wanda} and \emph{SparseGPT} \cite{frantar2023sparsegpt} leverage both weight magnitude and input activation norms (or inverse Hessian information) to prune weights efficiently without extensive retraining. While highly effective, these methods focus on pruning a single pre-trained model and do not address the complex spatial and representational conflicts that arise during multi-task expert merging.

\subsection{Test-Time Adaptation and Unsupervised Calibration}
Test-Time Adaptation (TTA) aims to adjust a pre-trained model to a target distribution during inference using only unlabeled test data. \emph{Tent} \cite{wang2020tent} pioneered TTA by minimizing the Shannon entropy of model predictions on-the-fly, optimizing batch normalization parameters to correct covariate shift. Similarly, \emph{MEMO} \cite{zhang2022memo} minimizes the entropy of a single test image augmented multiple times to enhance robustness to real-world corruptions. 

AdaMerging \cite{yang2024adamerging} successfully adopted this unsupervised entropy minimization loss to optimize merging coefficients without requiring ground-truth labels. However, optimizing parameters at test-time under non-differentiable operators (such as threshold-based magnitude pruning) is non-trivial. While Straight-Through Estimators (STE) \cite{bengio2013estimating} and zero-order search methods like Evolution Strategies (ES) \cite{salimans2017evolution} have been individually utilized to bypass non-differentiable thresholds in quantization-aware training or reinforcement learning, their application to co-optimizing sparse on-the-fly model merging is highly novel and addresses a significant practical bottleneck.
